\documentclass[a4paper, 11pt]{article}

\usepackage[utf8]{inputenc}
\RequirePackage[left=2cm,text={17cm, 24cm},top=3cm]{geometry}
\usepackage{times}
\usepackage{graphicx}
\usepackage{setspace}
\usepackage{hyperref}
\usepackage{fancyhdr}
\usepackage[outputdir=build]{minted}
\usepackage{amsmath}
\usepackage{textcomp}

% for indented itemize list
\usepackage{enumitem}

\graphicspath{{../src/outputs/}}

\begin{document}

    % header
    \begin{center}
    {\includegraphics[width=0.4\linewidth]{fig/logo_en.png}}
        \\
        [0.4cm]

        {\LARGE
        ISS - Project \\
        [0.4cm]
        }

        {\large
        Adrián Kálazi (xkalaz00) \\
        \today
        }
    \end{center}

    % actual content
    \section*{Abstract}

    The calculations were implemented using \textbf{Python (3.5+)} with the following libraries:

    \begin{itemize}
        \item numpy - arrays, \texttt{fft}, \texttt{ifft}
        \item scipy - \texttt{lfilter}
        \item soundfile - wav reading/writing
        \item matplotlib - plotting
    \end{itemize}

    The implementation consists of a Python module and a launcher script (\texttt{main.py}).

    Running the script:
    \begin{itemize}
        \item
        On a system without the required libraries - execute \texttt{./run-venv.sh} in the \texttt{src/} folder.

        This will setup a virtual environment, install the requirements and run the script.
        \item
        On a system with the required libraries - execute \texttt{python3 ./main.py} in the \texttt{src/} folder.
    \end{itemize}

    \section*{Solution}

    \begin{enumerate}
        \item
        The audio was recorded with \texttt{\href{https://apps.kde.org/en/kwave}{KWave}} in 16-bit resolution at 48~kHz.

        The sampling rate was later scaled down to 16~kHz using \texttt{ffmpeg}

        Durations of recorded signals: \\
        \begin{tabular}{ | l | l | l | }
            \hline
            \textbf{File}     & \textbf{Samples} & \textbf{Seconds} \\ \hline
            maskoff\_tone.wav & 24806            & 1.55             \\ \hline
            maskon\_tone.wav  & 17137            & 1.07             \\ \hline
        \end{tabular}

        \item
        The sentence recording process was the same as above

        Durations of recorded signals: \\
        \begin{tabular}{ | l | l | l | }
            \hline
            \textbf{File}         & \textbf{Samples} & \textbf{Seconds} \\ \hline
            maskoff\_sentence.wav & 41162            & 2.57             \\ \hline
            maskon\_sentence.wav  & 43312            & 2.71             \\ \hline
        \end{tabular}

        \newpage
        \item
        Frame duration $t_F = 20\,ms$ \\
        Sampling frequency $F_s = 16\,kHz$

        Frame size in samples: $n_F = t_F \times F_s = 20~ms \times 16~kHz = 320~samples$

        \includegraphics[width=\linewidth]{3_frames.pdf}

        \item
        Frames

        \includegraphics[width=\linewidth]{4_frame.pdf} \\
        \includegraphics[width=\linewidth]{4_frame_clipped.pdf} \\
        \includegraphics[width=\linewidth]{4_frame_autocorrelated.pdf} \\
        \includegraphics[width=\linewidth]{4_base_frequencies.pdf}

        \newpage
        Mean and variance of base frequencies: \\
        \begin{tabular}{ | l | l | l | }
            \hline
            \textbf{Signal}       & \textbf{Mean} $[Hz]$ & \textbf{Variance} $[Hz]$ \\ \hline
            maskoff\_sentence.wav & 139.466              & 0.803                    \\ \hline
            maskon\_sentence.wav  & 143.478              & 1.883                    \\ \hline
        \end{tabular}

        \vspace{2mm}
        There are two possible solutions that come to my mind to the mentioned problem:

        \begin{itemize}
            \item
            In our implementation, lag was defined as \texttt{np.argmax(autocorrelation\_array)} and therefore it was an integer.

            An alternate implementation would be to approximate lag as a weighted arithmetic mean from the surrounding indices and their values
            which would make it a floating point number resulting in better precision.
            This method would eliminate rapid changes of $f_0$ in most cases.

            \item
            Autocorrelation with higher resolution - this would eliminate the rapid changes in $f_0$
            but would require a longer computation time.
        \end{itemize}

        \item
        Spectrograms

        \includegraphics[width=\linewidth]{5_spectrum_maskoff.pdf} \\
        \includegraphics[width=\linewidth]{5_spectrum_maskon.pdf}

        \newpage

        DFT function implementation

        % @formatter:off
        % implementation from https://www.fit.vutbr.cz/study/courses/ISS/public/NEW_PRED/02_spectrum/spectrum.pdf slide 46
        \begin{minted}[resetmargins=true,tabsize=4, python3=true]{python}
        def dft(x):
            X = []
            N = len(x)
            for k in range(N):
                a = 0
                for n in range(N):
                    a += x[n] * cmath.exp(-2j * cmath.pi * (k/N) * n)
                X.append(a)
            return X
        \end{minted}
        % @formatter:on

        \item
        Frequency response

        Calculate the the mean from all frames, for each frequency

        \begin{gather*}
            H(e^{j\omega}) = \frac{1}{N} \sum_{n=0}^{N-1} \frac{ B_n(e^{j\omega}) }{ A_n(e^{j\omega}) } \\
            N \text{ - frame count, } A_n, B_n \text{ - spectrums for frame n}
        \end{gather*}

        \includegraphics[width=\linewidth]{6_frequency_response.pdf}

        \item
        Impulse response

        \includegraphics[width=\linewidth]{7_impulse_response.pdf}

        IDFT function implementation

        % @formatter:off
        % implementation from https://www.fit.vutbr.cz/study/courses/ISS/public/NEW_PRED/02_spectrum/spectrum.pdf slide 51
        \begin{minted}[resetmargins=true,tabsize=4, python3=true]{python}
        def idft(X):
            x = []
            N = len(X)
            for n in range(N):
                a = 0
                for k in range(N):
                    a += X[k] * cmath.exp(2j * cmath.pi * (k/N) * n)
                x.append(a / N)
            return x
        \end{minted}
        % @formatter:on

        \newpage

        \item
        Filtered signal

        \includegraphics[width=\linewidth]{8_signal_maskoff.pdf} \\
        \includegraphics[width=\linewidth]{8_signal_maskon.pdf} \\
        \includegraphics[width=\linewidth]{8_signal_filtered.pdf}

        \item
        Conclusion

    \end{enumerate}

    \newpage
    \section*{Solution - additional tasks}

    \begin{enumerate}
        \item[11.]
        Window function

        The Hann window function has been used.

        \includegraphics[width=\linewidth]{11_window.pdf} \\
        \includegraphics[width=\linewidth]{11_response.pdf}

        Phase shift (from task 15) has been applied before the window function,
        this choice was made because it shows better how the function behaves in the next graphs.

        \includegraphics[width=\linewidth]{11_frames_before.pdf} \\
        \includegraphics[width=\linewidth]{11_frames_after.pdf}

        \newpage
        Frame spectral densities

        \includegraphics[width=\linewidth]{11_frame_spectrum_before.pdf} \\
        \includegraphics[width=\linewidth]{11_frame_spectrum_after.pdf}

        From the above graphs we can see that the window function
        supresses the frequencies on the edge of the spectrum and the resulting frequency spectrum is smoother

        \newpage
        \item[15.]
        Phase shift

        \includegraphics[width=\linewidth]{15_frames_before.pdf}
        \includegraphics[width=\linewidth]{15_frames_after.pdf}

        \includegraphics[width=\linewidth]{15_shifts.pdf}

        From the above graph, we can see that the required phase shift increases almost linearly -
        this is caused by the frequency of the "mask off" signal being a little bit lower than the "mask on" signal

        \vspace{6mm}

        \textbf{Q:} When we take the sum of phase shifts from both sides,
        we can see that it resembles the lag in autocorrelation,
        what is the reason for this behaviour?

        \textbf{A:} This is due to the coefficients (and lag) representing a shift at which both signals are the "most similar".

        In autocorrelation this results in the lag representing the number of samples until the next period occurs.

        Because autocorrelation is done on only one signal,
        the shortest shift at which the signal is similar to itself is its own period.

        \textbf{The autocorrelation coefficient (and therefore lag) represents one period.}

        In cross-correlation the coefficient represents the shortest shift in a certain direction,
        but because the signals have similar frequencies and the cross-correlation is done from both sides,
        summing up these two values results in them
        being equal to the signals period $T$ in samples.

        \textbf{Both cross-correlation coefficients represent a part of the period (e.g.\ 1/3 and 2/3)
        which summed up results in one period.}
    \end{enumerate}

\end{document}
